%% kapa: exam=1 round=25 year=2014 period=1 form=A subject=civil_law subject_ko=민법 seq=5 number=5
%% source: Q-Net 감정평가사 기출문제 (cst00309) artlSeq=5101207 file=extracted/1008/5101207/2014년 제25회 감정평가사 1차시험 기출문제/1교시(A형).hwp
%% conversion: kordoc 4.24.51 -> tools/gen1.py (auto); review=proofread-sonnet; flags=none
\begin{question}{5}
미성년자에 관한 설명으로 옳지 않은 것은? (다툼이 있으면 판례에 의함)
\begin{choices}
  \item 1995. 3. 30. 오후 9시에 출생한 자는 2014. 3. 29. 오후 12시에 성년자로 된다.
  \item 법정대리인이 미성년자에게 특정한 영업을 허락한 경우에 미성년자는 그 영업에 관하여는 성년자와 동일한 행위능력을 가진다.
  \item 미성년자가 부담 없는 증여를 받는 행위는 법정대리인의 동의를 요하지 않는다.
  \item 미성년자가 월 소득범위 내에서 소규모의 일상적인 신용구매계약을 체결하였더라도 스스로 얻은 소득에 대해 법정대리인의 묵시적 처분허락이 있었다고 볼 수는 없다.
  \item 미성년자는 자신의 노무제공에 대한 임금의 청구를 독자적으로 할 수 있다.
\end{choices}
\end{question}
\answer{4}
